% Luc Destombe --- licence CC BY-NC-SA 4.0
% https://creativecommons.org/licenses/by-nc-sa/4.0/deed.fr
% Fichier autonome : compiler avec pdflatex, deux fois (signets, nombre de pages).
\def\lyceeversion{1.0}
\documentclass[11pt,a4paper]{article}
\makeatletter
\newif\ifeleve
\elevefalse

\newif\iflycee@palatino
\lycee@palatinotrue

\RequirePackage[utf8]{inputenc}
\RequirePackage[T1]{fontenc}
\RequirePackage[french]{babel}
\RequirePackage{etoolbox}

\newcommand{\ShorthandsOff}{\@ifundefined{shorthandoff}{}{\shorthandoff{;:!?}}}
\newcommand{\ShorthandsOn}{\@ifundefined{shorthandon}{}{\shorthandon{;:!?}}}
\ShorthandsOff
\AtBeginDocument{\ShorthandsOn}
\AtBeginEnvironment{tikzpicture}{\ShorthandsOff}

\RequirePackage{geometry}
\RequirePackage{fancyhdr}
\RequirePackage{lastpage}
\RequirePackage{multicol}
\RequirePackage{array}
\RequirePackage{enumitem}
\RequirePackage{xcolor}
\RequirePackage{needspace}

\RequirePackage{amsmath,amssymb}
\RequirePackage{mathpazo}
\renewcommand{\ttdefault}{lmtt}
\DeclareMathSizes{10.95}{10.95}{8}{5.5}
\begingroup\catcode`\_=\active \catcode`\^=\active
\gdef_#1{\sb{\scriptscriptstyle #1}}
\gdef^#1{\sp{\scriptscriptstyle\lycee@moinsepais #1}}
\endgroup
\newcommand\lycee@moins{\mathbin{%
  \setbox\z@\hbox{$\m@th\scriptscriptstyle\lycee@moinsorig$}%
  \dimen@=.55\wd\z@ \advance\dimen@-.5pt
  \hbox to.75\wd\z@{\hss$\m@th\scriptscriptstyle\vcenter{\hbox{%
    \tikz\draw[line width=.5pt,line cap=round](0,0)--(\the\dimen@,0);}}$\hss}}}
\begingroup\lccode`\~=`\- \lowercase{\endgroup\let~}\lycee@moins
\newcommand\lycee@moinsepais{\mathcode`\-="8000 }
\AtBeginDocument{\mathchardef\lycee@moinsorig=\mathcode`\-\relax}
\AtBeginDocument{\catcode`\_=12 \mathcode`\_="8000
  \catcode`\^=12 \mathcode`\^="8000 }
\newcommand\lycee@exposants{%
  \fontdimen13\textfont2=.48\fontdimen6\textfont2
  \fontdimen14\textfont2=.48\fontdimen6\textfont2
  \fontdimen15\textfont2=.48\fontdimen6\textfont2
  \fontdimen13\scriptfont2=.48\fontdimen6\scriptfont2
  \fontdimen14\scriptfont2=.48\fontdimen6\scriptfont2
  \fontdimen15\scriptfont2=.48\fontdimen6\scriptfont2}
\every@math@size\expandafter{\the\every@math@size\lycee@exposants}
\newlength{\retraitcalculs}
\setlength{\retraitcalculs}{2em}

\RequirePackage{tikz}
\usetikzlibrary{arrows.meta,calc,babel,shapes.misc}
\RequirePackage[most]{tcolorbox}
\tcbuselibrary{skins,breakable}

\definecolor{coulDef}{HTML}{1F4E79}
\definecolor{coulProp}{HTML}{7030A0}
\definecolor{coulRem}{HTML}{C55A11}
\definecolor{coulEx}{HTML}{2E7D32}
\definecolor{coulExo}{HTML}{B03A2E}
\definecolor{coulPart}{HTML}{1F4E79}
\definecolor{coulMeth}{HTML}{1B6E4F}
\definecolor{coulCourbe}{HTML}{0B5ED7}
\definecolor{coulCourbeB}{HTML}{C00000}
\definecolor{coulCourbeC}{HTML}{2E7D32}
\definecolor{coulGrille}{HTML}{8C8C8C}
\definecolor{coulRep}{HTML}{1B6E4F}
\definecolor{coulBar}{HTML}{2E7D32}

\newcommand{\R}{\mathbb{R}}
\newcommand{\Rs}{\mathbb{R}^{*}}
\renewcommand{\le}{\leqslant}
\renewcommand{\ge}{\geqslant}
\renewcommand{\approx}{\simeq}
\newcommand{\vect}[1]{\overrightarrow{#1}}
\newcommand{\norme}[1]{\left\lVert #1 \right\rVert}
\newcommand{\intervalle}[2]{\left[#1\,;#2\right]}
\RequirePackage{cancel}
\renewcommand{\CancelColor}{\color{coulExo}}
\newcommand{\fois}[1]{\times{\color{coulExo}#1}}
\newcommand{\barre}[1]{\cancel{#1}}
\newlength{\lycee@eqh}\newlength{\lycee@eqd}
\newsavebox{\lycee@eqa}\newsavebox{\lycee@eqb}\newsavebox{\lycee@eqc}
\newcommand{\lycee@eqcol}[1]{\hbox{$\renewcommand{\arraystretch}{1.35}%
  \begin{array}[t]{@{}>{\displaystyle}r@{}>{\displaystyle{}}l@{}}#1\end{array}$}}
\newcommand{\lycee@eqmesure}[1]{%
  \ifdim\ht#1>\lycee@eqh\setlength{\lycee@eqh}{\ht#1}\fi
  \ifdim\dp#1>\lycee@eqd\setlength{\lycee@eqd}{\dp#1}\fi}
\newcommand{\lycee@eqtrait}{\hspace{1.2em}\tikz[baseline=0pt]\draw[dash pattern=on 2pt off 2pt,
  line width=0.5pt] (0,\lycee@eqh) -- (0,-\lycee@eqd);\hspace{1.2em}}
\newcommand{\eqproduit}[2]{\par\smallskip\noindent C'est une équation produit :
  $A\times B=0 \iff A=0$ ou $B=0$.\par\nopagebreak\smallskip
  \begingroup\setlength{\lycee@eqh}{0pt}\setlength{\lycee@eqd}{0pt}%
  \setbox\lycee@eqa=\lycee@eqcol{#1}\setbox\lycee@eqb=\lycee@eqcol{#2}%
  \lycee@eqmesure\lycee@eqa\lycee@eqmesure\lycee@eqb%
  \noindent\hspace*{\retraitcalculs}\usebox\lycee@eqa\lycee@eqtrait\usebox\lycee@eqb\endgroup\par\smallskip}
\newcommand{\eqproduittrois}[3]{\par\smallskip\noindent C'est une équation produit :
  $A\times B\times C=0 \iff A=0$ ou $B=0$ ou $C=0$.\par\nopagebreak\smallskip
  \begingroup\setlength{\lycee@eqh}{0pt}\setlength{\lycee@eqd}{0pt}%
  \setbox\lycee@eqa=\lycee@eqcol{#1}\setbox\lycee@eqb=\lycee@eqcol{#2}\setbox\lycee@eqc=\lycee@eqcol{#3}%
  \lycee@eqmesure\lycee@eqa\lycee@eqmesure\lycee@eqb\lycee@eqmesure\lycee@eqc%
  \noindent\hspace*{\retraitcalculs}\usebox\lycee@eqa\lycee@eqtrait\usebox\lycee@eqb\lycee@eqtrait\usebox\lycee@eqc\endgroup
  \par\smallskip}

\tikzset{grille/.style={coulGrille,line width=0.4pt}}
\newcommand{\quadrillage}[4]{\draw[grille,step=1] (#1,#2) grid (#3,#4);}
\tikzset{
  croix/.style={cross out,draw,line width=0.6pt,inner sep=0pt,minimum size=4pt},
  croixdroite/.style={croix,rotate=45,line width=0.8pt},
}
\newcommand{\pt}[3]{\path (#1) node[croix]{} node[#2,font=\small,inner sep=2.5pt]{$#3$};}

\newcommand{\lycee@repere}[4]{%
  \draw[grille,step=1] (#1,#2) grid (#3,#4);
  \draw[-{Stealth[length=2mm]},thick] (#1,0)--({#3+0.4},0) node[below left=-1pt]{$x$};
  \draw[-{Stealth[length=2mm]},thick] (0,#2)--(0,{#4+0.4}) node[below left=-1pt]{$y$};
  \node[below left,font=\scriptsize,inner sep=1.5pt] at (0,0) {$0$};}
\newcommand{\lycee@gradx}[1]{\foreach \k/\l in {#1}{%
  \draw (\k,0.07)--(\k,-0.07) node[below,font=\scriptsize,inner sep=1.5pt]{$\l$};}}
\newcommand{\lycee@grady}[1]{\foreach \k/\l in {#1}{%
  \draw (0.07,\k)--(-0.07,\k) node[left,font=\scriptsize,inner sep=1.5pt]{$\l$};}}
\newcommand{\lycee@intff}[2]{\left[#1\,;#2\right]}
\newcommand{\lycee@intfo}[2]{\left[#1\,;#2\right[}
\newcommand{\lycee@intof}[2]{\left]#1\,;#2\right]}
\newcommand{\lycee@intoo}[2]{\left]#1\,;#2\right[}
\newcommand{\lycee@ens}[1]{\left\{#1\right\}}
\AtBeginDocument{%
  \providecommand{\repere}{\lycee@repere}%
  \providecommand{\gradx}{\lycee@gradx}%
  \providecommand{\grady}{\lycee@grady}%
  \providecommand{\Cf}{\mathcal{C}\sb{\scriptscriptstyle f}}%
  \providecommand{\Cg}{\mathcal{C}\sb{\scriptscriptstyle g}}%
  \providecommand{\intff}{\lycee@intff}%
  \providecommand{\intfo}{\lycee@intfo}%
  \providecommand{\intof}{\lycee@intof}%
  \providecommand{\intoo}{\lycee@intoo}%
  \providecommand{\ens}{\lycee@ens}}

\newlist{questions}{enumerate}{3}
\setlist[questions,1]{label=\arabic*\textdegree),leftmargin=*,itemsep=2pt,topsep=3pt}
\setlist[questions,2]{label=\alph*),leftmargin=*,itemsep=1pt,topsep=2pt}
\setlist[questions,3]{label=\roman*),leftmargin=*,itemsep=1pt,topsep=1pt}
\setlist[itemize]{label=\textbullet,leftmargin=*,itemsep=2pt,topsep=3pt}

\newcommand{\besoinplace}[1]{\par\penalty\@M\begingroup
  \dimen@=#1\relax\dimen@ii=\pagegoal\advance\dimen@ii-\pagetotal
  \ifdim\dimen@>\dimen@ii\ifdim\dimen@ii>\z@\vfil\fi\break\fi\endgroup}

\newcommand{\lycee@pied}{\thepage/\pageref{LastPage}}
\newcommand{\entete}[2]{%
  \pagestyle{fancy}\fancyhf{}%
  \fancyhead[L]{\footnotesize\color{coulPart}\textbf{#1}}%
  \fancyhead[R]{\footnotesize\color{coulPart}\textbf{#2}}%
  \fancyfoot[C]{\footnotesize\lycee@pied}%
  \renewcommand{\headrulewidth}{0.4pt}%
  \renewcommand{\headrule}{\hbox to\headwidth{%
    \color{coulPart!40}\leaders\hrule height \headrulewidth\hfill}}}

\RequirePackage{ccicons}
\newcommand{\lycee@licence}{{\scriptsize\color{gray}%
  \copyright\ Luc Destombe --- \ccbyncsa\ CC BY-NC-SA 4.0}}
\newif\iflycee@licence \lycee@licencetrue
\newcommand{\sanslicence}{\lycee@licencefalse}
\AtBeginDocument{\iflycee@licence\AddToHookNext{shipout/foreground}{%
  \put(0,0){\rlap{\kern\dimexpr1in+\hoffset+\oddsidemargin+\textwidth\relax
    \llap{\raisebox{-\dimexpr1in+\voffset+\topmargin+\headheight+\headsep
      +\textheight+\footskip\relax}{\lycee@licence}}}}}\fi}

\newcommand{\lycee@site}{faire-des-maths.fr}
\newcommand{\lycee@version}{{\scriptsize\color{gray}\lycee@site
  \ifdefined\lyceeversion\ --- Version~\lyceeversion\fi}}
\AtBeginDocument{\iflycee@licence\AddToHookNext{shipout/foreground}{%
  \put(0,0){\rlap{\kern\dimexpr1in+\hoffset+\oddsidemargin\relax
    \raisebox{-\dimexpr1in+\voffset+\topmargin+\headheight+\headsep
      +\textheight+\footskip\relax}{\lycee@version}}}}\fi}

\newcommand{\formulesserrees}{%
  \setlength{\abovedisplayskip}{3pt}\setlength{\belowdisplayskip}{3pt}%
  \setlength{\abovedisplayshortskip}{2pt}\setlength{\belowdisplayshortskip}{3pt}}

\setlength{\parindent}{0pt}

\RequirePackage[implicit=false,hidelinks,pdfencoding=auto,psdextra,bookmarksopen,
  bookmarksopenlevel=1,pdfcreator={},pdfproducer={}]{hyperref}
\RequirePackage{bookmark}
\hypersetup{pdfauthor={Luc Destombe},pdfkeywords={CC BY-NC-SA 4.0}}
\newcounter{lycee@signet}
\newcommand{\signet}[3][]{\stepcounter{lycee@signet}%
  \edef\lycee@dest{\if\relax\detokenize{#1}\relax
    signet.\arabic{lycee@signet}\else#1\fi}%
  \hypertarget{\lycee@dest}{}\bookmark[level=#2,dest=\lycee@dest]{#3}}
\pdfstringdefDisableCommands{%
  \def\R{R}\def\vect#1{#1}\def\Cf{Cf}\def\Cg{Cg}\def\fois#1{×#1}%
  \def\barre#1{#1}\def\intervalle#1#2{[#1;#2]}\def\intff#1#2{[#1;#2]}%
  \def\textsuperscript#1{#1}\def\dfrac#1#2{#1/#2}\def\frac#1#2{#1/#2}%
  \def\vec#1{#1⃗}}%
\newcommand{\res}[1]{#1}
\geometry{top=2cm,bottom=1cm,left=1cm,right=1cm,headheight=14pt,footskip=0.6cm}

\RequirePackage{pgfplots}
\pgfplotsset{compat=1.18}
\RequirePackage{tkz-tab}

\newcounter{partiecours}
\renewcommand{\thepartiecours}{\Roman{partiecours}}
\newcommand{\partiecours}[1]{%
  \refstepcounter{partiecours}%
  \Needspace*{8\baselineskip}%
  \bigskip
  \noindent\begin{tcolorbox}[enhanced,colback=coulPart,colframe=coulPart,
      boxrule=0pt,arc=2pt,left=8pt,right=8pt,top=4pt,bottom=4pt,
      before skip=6pt,after skip=10pt]
    \color{white}\bfseries\large \thepartiecours{} --- #1%
    \signet[partie-\arabic{partiecours}]{1}{\thepartiecours{} --- #1}
  \end{tcolorbox}%
  \addcontentsline{toc}{section}{\thepartiecours{} --- #1}%
}

\newtcolorbox{boitecours}[3][]{%
  enhanced,breakable,
  colback=white,colframe=#2,coltitle=white,
  fonttitle=\bfseries,
  attach boxed title to top left={xshift=8pt,yshift=-\tcboxedtitleheight/2},
  boxed title style={colback=#2,arc=2pt,boxrule=0pt},
  boxrule=0.9pt,arc=2pt,
  left=8pt,right=8pt,top=8pt,bottom=6pt,
  title={#3},#1}

\newcounter{methode}
\newenvironment{definitionbox}[1][Définition]%
  {\begin{boitecours}[phantom={\signet{2}{#1}}]{coulDef}{#1}}{\end{boitecours}}
\newenvironment{proprietebox}[1][Propriété]%
  {\begin{boitecours}[phantom={\signet{2}{#1}}]{coulProp}{#1}}{\end{boitecours}}
\newenvironment{methodebox}[1][Méthode]{\stepcounter{methode}%
  \begin{boitecours}[phantom={\signet[methode-\arabic{methode}]{2}{#1}}]{coulMeth}{#1}}%
  {\end{boitecours}}

\newcommand{\etape}[1]{\par\smallskip\textbf{\color{coulMeth}#1}\par\nobreak\smallskip}
\newcommand{\enonce}[1]{\emph{#1}\par\smallskip}
\newcommand{\reponse}[1]{\par\smallskip
  {\footnotesize\itshape\color{black!55}Réponses : #1}}
\newcommand{\demo}[1]{\textbf{\color{coulMeth}Démonstration.} #1}
\newcommand{\afaire}[1]{%
  \par\smallskip
  \noindent{\small\color{coulExo}$\blacktriangleright$\ \textbf{À faire :}
    fiche d'exercices du chapitre, #1.}\par\smallskip}

\newenvironment{calculs}{\setlength{\jot}{5pt}\@fleqntrue\@mathmargin\retraitcalculs
  \setlength{\abovedisplayskip}{4pt}\setlength{\belowdisplayskip}{4pt}%
  \setlength{\abovedisplayshortskip}{4pt}\setlength{\belowdisplayshortskip}{4pt}%
  \csname alignat*\endcsname{2}}%
  {\csname endalignat*\endcsname}
\newcommand{\rem}[1]{\text{\small\itshape\color{black!60}#1}}
\newcommand{\explique}[2]{\par\smallskip\noindent
  \begin{minipage}[t]{0.57\linewidth}\raggedright #1\end{minipage}\hfill
  \begin{minipage}[t]{0.40\linewidth}\raggedright\leavevmode
    {\small\itshape\color{black!60}#2}\end{minipage}\par}
\newcommand{\cas}[1]{\par\smallskip\noindent\textbf{\color{coulExo}#1}\vadjust{\nobreak}}
\newcommand{\calc}[1]{\par\smallskip\textbf{\color{coulExo}#1}\par\nobreak}

\newtcolorbox{remarquebox}[1][Remarque]{%
  enhanced,breakable,colback=white,colframe=coulRem,
  boxrule=0pt,leftrule=2.5pt,arc=0pt,outer arc=0pt,
  left=8pt,right=6pt,top=5pt,bottom=5pt,
  before upper={\textbf{\color{coulRem}#1}\par\smallskip}}

\newtcolorbox{exemplebox}[1][Exemple]{%
  enhanced,breakable,colback=white,colframe=coulEx,
  boxrule=0pt,leftrule=2.5pt,arc=0pt,outer arc=0pt,
  left=8pt,right=6pt,top=5pt,bottom=5pt,
  before upper={\textbf{\color{coulEx}#1}\par\smallskip}}

\pgfplotsset{
  repere/.style={
    axis lines=middle,
    axis line style={-{Stealth[length=2.4mm]},thick},
    every axis x label/.style={at={(ticklabel* cs:1.0)},anchor=west},
    every axis y label/.style={at={(ticklabel* cs:1.0)},anchor=south},
    label style={font=\footnotesize},
    tick label style={font=\scriptsize},
    grid=both,
    major grid style={grille},
    minor tick num=0,
    tick align=inside,
    clip mode=individual,
  },
}
\tikzset{
  courbe/.style={coulCourbe,line width=1.1pt,smooth},
  courbeB/.style={coulCourbeB,line width=1.1pt,smooth},
  courbeC/.style={coulCourbeC,line width=1.1pt,smooth},
}

\newlist{etapes}{enumerate}{1}
\setlist[etapes]{label=\textbf{\arabic*.},leftmargin=*,itemsep=1pt,topsep=2pt}

\newlist{expressions}{enumerate}{1}
\setlist[expressions]{label={\Alph*\ $=$},leftmargin=*,itemsep=3pt,topsep=3pt}

\newcounter{exo}
\renewcommand{\theexo}{\ifnum\value{exo}<10 0\fi\arabic{exo}}
\newtcolorbox{exercice}[1][]{%
  enhanced,breakable,
  colback=white,colframe=coulExo,
  boxrule=0.9pt,arc=2pt,
  left=8pt,right=8pt,top=8pt,bottom=6pt,
  coltitle=white,fonttitle=\bfseries,
  attach boxed title to top left={xshift=8pt,yshift=-\tcboxedtitleheight/2},
  boxed title style={colback=coulExo,arc=2pt,boxrule=0pt},
  phantom={\signet{2}{Exercice \arabic{exo}}},
  title={Exercice \theexo\ifx\relax#1\relax\else{} --- #1\fi}}
\newenvironment{exo}[1][\relax]{\refstepcounter{exo}\begin{exercice}[#1]}{\end{exercice}}

  \newenvironment{correction}{\par}{\par}
  \newcommand{\autableau}{}
  \newcommand{\etapeexemples}{\etape{Exemples corrigés}}
  \newcommand{\etapetableau}[1]{\etape{#1}}
  \newenvironment{exempletableau}[1][Exemple]%
    {\begin{exemplebox}[#1]}{\end{exemplebox}}

\RequirePackage{comment}
\excludecomment{correctionavj}

\newcommand{\coursEntete}{}
\newcommand{\coursNiveau}{}
\newcommand{\coursTitre}{}
\newcommand{\coursSurTitre}{}
\newcommand{\coursSousTitre}{}

\newcommand{\enteteCours}[1]{\renewcommand{\coursEntete}{#1}}
\newcommand{\chapitrecours}[2]{\enteteCours{Chapitre #1 --- #2}}
\newcommand{\niveaucours}[1]{\renewcommand{\coursNiveau}{#1}}
\newcommand{\titrecours}[3][]{%
  \renewcommand{\coursTitre}{#2}%
  \renewcommand{\coursSurTitre}{#1}%
  \renewcommand{\coursSousTitre}{#3}}

\pagestyle{fancy}
\fancyhf{}
\fancyhead[L]{\footnotesize\color{coulPart}\textbf{\coursEntete}}
\fancyhead[R]{\footnotesize\color{coulPart}\textbf{\coursNiveau}}
\fancyfoot[C]{\footnotesize\thepage}
\renewcommand{\headrulewidth}{0.4pt}
\renewcommand{\headrule}{\hbox to\headwidth{\color{coulPart!40}\leaders\hrule height \headrulewidth\hfill}}

\setlength{\parskip}{2pt}

\newcommand{\lycee@mentioneleve}{\\[2pt]{\small\itshape Version élève}}
\AtBeginDocument{%
  \begin{center}
    \begin{tcolorbox}[enhanced,width=0.72\linewidth,colback=white,colframe=coulPart,
        boxrule=1.6pt,arc=3pt,halign=center,top=10pt,bottom=10pt]
      \ifx\coursSurTitre\empty
        {\Huge\bfseries\color{coulPart} \coursTitre}\\[4pt]
      \else
        {\Huge\bfseries\color{coulPart} \coursTitre}\\[3pt]
        {\Large\bfseries\color{coulPart} \coursSurTitre}\\[5pt]
      \fi
      {\normalsize \coursSousTitre}%
      
    \end{tcolorbox}
  \end{center}%
}

\RequirePackage{tikz}
\RequirePackage{xcolor}

\usetikzlibrary{decorations.text}
\definecolor{lycee@metal}{HTML}{1F4E79}
\definecolor{lycee@manche}{HTML}{8A5A2B}
\definecolor{lycee@etincelle}{HTML}{F28C28}
\newcommand{\lycee@pluriel}[2]{#1 #2\ifnum#1>1 s\fi}
\newcommand{\lycee@sceau}{%
  \edef\lycee@bas{\lycee@pluriel{\lyceerelectures}{relecture} et
    \lycee@pluriel{\lyceerodages}{rodage}}%
  \begin{tikzpicture}[line join=round]
    \fill[coulPart!10,draw=coulPart,line width=.8pt] (0:2.2)
      \foreach \i in {1,...,71}{ -- (\i*5:{isodd(\i)?2.1:2.2}) } -- cycle;
    \draw[coulPart,line width=1.2pt] (0,0) circle (1.93);
    \draw[coulPart,line width=.4pt] (0,0) circle (1.85);
    \draw[coulPart,line width=.4pt] (0,0) circle (1.2);
    \path[decorate,decoration={text along path,text align=center,
        text={|\fontsize{9.5}{11}\selectfont\bfseries\color{coulPart}|TEST{\'E} AVEC MES {\'E}L{\`E}VES}}]
      (180:1.45) arc (180:0:1.45);
    \path[decorate,decoration={text along path,text align=center,
        text/.expanded={|\noexpand\fontsize{8.5}{10}\noexpand\selectfont\noexpand\color{coulPart}|\lycee@bas}}]
      (195:1.68) arc (195:345:1.68);
    \begin{scope}[x=.25mm,y=-.25mm]
      \begin{scope}[shift={(-32,-40)}]%
        \fill[lycee@metal] (23,31) rectangle (41,49);
        \draw[lycee@etincelle,line width=.5pt,line cap=round]
          (20,30)--(15,27) (19,34)--(14,35) (44,30)--(49,27) (45,34)--(50,35);
        \fill[lycee@manche,rounded corners=.3mm] (38,24.6) rectangle (60,27.6);
        \fill[lycee@metal,rounded corners=.4mm] (25,20) rectangle (39,31);
        \fill[lycee@etincelle,rounded corners=.2mm] (25,20) rectangle (28,31);
      \end{scope}
    \end{scope}
    \ifdefined\lyceesceauplusieurs\def\lycee@annee{Ann{\'e}es~:}\else
      \def\lycee@annee{Ann{\'e}e}\fi
    \path[decorate,decoration={text along path,text align=center,
        text/.expanded={|\noexpand\fontsize{8}{10}\noexpand\selectfont\noexpand\color{coulPart}|%
          \unexpanded\expandafter{\lycee@annee} \unexpanded\expandafter{\lyceesceauannees}}}]
      (95:1.08) arc (95:445:1.08);
  \end{tikzpicture}}
\AtEndDocument{\ifdefined\lyceesceauannees
  \par\besoinplace{5cm}\vspace{1.5em}%
  \begin{center}\lycee@sceau\end{center}\fi}
\makeatother

\chapitrecours{12}{Statistiques}
\niveaucours{Seconde}
\titrecours{STATISTIQUES}{Cours --- Seconde}

\begin{document}

\tikzset{
  barre/.style={fill=coulCourbe!20,draw=coulCourbe,line width=0.8pt},
  secteur/.style={draw=black,line width=0.6pt},
}

\partiecours{Effectifs, fréquences, effectifs cumulés, classes}

\begin{definitionbox}[Définition 1 --- Vocabulaire]
\begin{itemize}
  \item La \textbf{population} est l'ensemble étudié (les élèves d'une classe, les
    ménages français) ; chacun de ses membres est un \textbf{individu}.
  \item Le \textbf{caractère} est ce que l'on observe sur chaque individu. Il est
    \textbf{qualitatif} quand ses valeurs ne sont pas des nombres (moyen de transport,
    matière préférée) ; \textbf{quantitatif} sinon, \textbf{discret} s'il ne prend
    que des valeurs isolées (nombre de livres lus), \textbf{continu} s'il peut prendre
    toutes les valeurs d'un intervalle (durée d'un trajet).
  \item La liste des valeurs relevées est une \textbf{série statistique}. Les
    réponses « ne se prononce pas » n'en font pas partie.
\end{itemize}
\end{definitionbox}

\begin{definitionbox}[Définition 2 --- Effectif, fréquence]
\begin{itemize}
  \item L'\textbf{effectif} d'une valeur est le nombre d'individus qui ont cette
    valeur ; l'\textbf{effectif total} $N$ est le nombre de valeurs de la série.
  \item La \textbf{fréquence} d'une valeur est la proportion des individus qui ont
    cette valeur (chapitre 7) :
    \[
    f=\frac{\text{effectif de la valeur}}{\text{effectif total}}
    \]
    C'est un nombre entre $0$ et $1$, écrit en fraction, en écriture décimale ou en
    pourcentage. La somme des fréquences vaut $1$, soit $100\,\%$.
\end{itemize}
\end{definitionbox}

\begin{definitionbox}[Définition 3 --- Effectifs et fréquences cumulés]
Les valeurs étant rangées dans l'ordre croissant, l'\textbf{effectif cumulé croissant}
d'une valeur est la somme des effectifs des valeurs qui lui sont inférieures ou
égales. Il répond à la question « combien d'individus ont au plus cette valeur ? ».
La \textbf{fréquence cumulée croissante} se définit de même avec les fréquences.
\end{definitionbox}

\begin{methodebox}[Méthode 1 --- Dresser un tableau d'effectifs, de fréquences et d'effectifs cumulés]
\etapeexemples
Dans une classe de seconde, $30$ élèves ont répondu à la question « Combien de
séances de sport faites-vous par semaine, en dehors de l'EPS ? ». Deux ont répondu
« ne se prononce pas » ; voici les réponses des $28$ autres.
\begin{center}
$0$ ; $2$ ; $1$ ; $3$ ; $2$ ; $0$ ; $4$ ; $2$ ; $1$ ; $0$ ; $2$ ; $3$ ; $1$ ; $2$ ;
$0$ ; $4$ ; $2$ ; $1$ ; $3$ ; $0$ ; $2$ ; $1$ ; $0$ ; $4$ ; $2$ ; $3$ ; $1$ ; $0$
\end{center}
\begin{questions}[label=\alph*)]
  \item Préciser la population, le caractère étudié et son type.
  \item Dresser le tableau des effectifs, des fréquences (arrondies au centième) et
    des effectifs cumulés croissants.
  \item Combien d'élèves font au plus $2$ séances ? Quelle proportion des élèves
    cela représente-t-il ?
\end{questions}
\begin{correction}
\cas{a)}
\explique{Population : les \res{$28$ élèves qui ont répondu}.}{les « ne se prononce
pas » sont retirés}
\explique{Caractère : le nombre de séances de sport par semaine, \res{quantitatif
discret}.}{des nombres entiers, isolés}
\cas{b)}
\explique{On compte chaque valeur : sept fois $0$, six fois $1$, huit fois $2$, quatre
fois $3$ et trois fois $4$.}{on barre chaque valeur comptée}
\begin{calculs}
N &= 7+6+8+4+3 &\qquad& \rem{somme des effectifs}\\
N &= \res{28} && \rem{on retrouve les $28$ réponses}
\end{calculs}
\begin{calculs}
f &= \dfrac{7}{28} &\qquad& \rem{fréquence de la valeur $0$}\\
f &= \dfrac{1}{4} &&\\
f &= \res{0{,}25} && \rem{soit $25\,\%$}
\end{calculs}
\explique{Les effectifs cumulés s'obtiennent en ajoutant chaque effectif au précédent :
$7$, puis $7+6=13$, puis $13+8=21$\ldots}{définition 3}
\begin{center}
\renewcommand{\arraystretch}{1.5}
\begin{tabular}{|l|*{5}{c|}c|}
\hline
Nombre de séances & $0$ & $1$ & $2$ & $3$ & $4$ & Total\\\hline
Effectif & $7$ & $6$ & $8$ & $4$ & $3$ & $28$\\\hline
Fréquence & $0{,}25$ & $0{,}21$ & $0{,}29$ & $0{,}14$ & $0{,}11$ & $1$\\\hline
Effectif cumulé croissant & $7$ & $13$ & $21$ & $25$ & $28$ & \\\hline
\end{tabular}
\end{center}
\cas{c)}
\explique{L'effectif cumulé de la valeur $2$ vaut $21$ : \res{$21$ élèves} font au plus
$2$ séances.}{au plus deux : zéro, une ou deux séances}
\begin{calculs}
f &= \dfrac{21}{28} &\qquad& \rem{fréquence cumulée de la valeur $2$}\\
f &= \dfrac{3}{4} && \rem{on simplifie par $7$}\\
f &= \res{0{,}75} && \rem{soit $75\,\%$ des élèves}
\end{calculs}
\end{correction}

\etape{À vous de jouer}
Dans une autre classe, $32$ élèves ont répondu à la question « Combien de livres
avez-vous lus depuis la rentrée, en dehors des lectures demandées en cours ? ». Deux
ont répondu « ne se prononce pas » ; voici les réponses des $30$ autres.
\begin{center}
$1$ ; $0$ ; $2$ ; $1$ ; $3$ ; $4$ ; $1$ ; $0$ ; $2$ ; $2$ ; $1$ ; $3$ ; $0$ ; $1$ ; $4$ ;
$2$ ; $1$ ; $0$ ; $3$ ; $2$ ; $1$ ; $1$ ; $0$ ; $2$ ; $3$ ; $4$ ; $1$ ; $2$ ; $0$ ; $3$
\end{center}
\begin{questions}[label=\alph*)]
  \item Préciser la population, le caractère étudié et son type.
  \item Dresser le tableau des effectifs, des fréquences (arrondies au centième) et
    des effectifs cumulés croissants.
  \item Combien d'élèves ont lu au plus $2$ livres ? Quelle proportion des élèves
    cela représente-t-il ?
\end{questions}
\reponse{a) les $30$ élèves qui ont répondu ; nombre de livres lus, quantitatif
discret ; b)~effectifs $6$ ; $9$ ; $7$ ; $5$ ; $3$ ; fréquences $0{,}2$ ; $0{,}3$ ;
$0{,}23$ ; $0{,}17$ ; $0{,}1$ ; effectifs cumulés $6$ ; $15$ ; $22$ ; $27$ ; $30$ ;
c)~$22$ élèves, soit environ $73\,\%$.}
\end{methodebox}

\afaire{exercices 1 à 4}

\begin{definitionbox}[Définition 4 --- Regroupement en classes]
Quand les valeurs sont nombreuses et toutes différentes (caractère continu), on les
regroupe en \textbf{classes} de même amplitude. La classe $\intfo{a}{b}$ contient les
valeurs $x$ telles que $a\le x<b$ ; son \textbf{amplitude} est $b-a$.
\enonce{Exemple : une durée de $20$ min est dans la classe $\intfo{20}{30}$, pas
dans $\intfo{10}{20}$.}
\end{definitionbox}

\begin{methodebox}[Méthode 2 --- Regrouper une série en classes]
\etapeexemples
Dans une classe de seconde, $24$ élèves ont indiqué leur temps de trajet entre leur
domicile et le lycée, en minutes.
\begin{center}
$5$ ; $12$ ; $25$ ; $8$ ; $18$ ; $35$ ; $15$ ; $22$ ; $10$ ; $42$ ; $28$ ; $14$ ;
$7$ ; $19$ ; $30$ ; $16$ ; $45$ ; $11$ ; $24$ ; $3$ ; $17$ ; $33$ ; $20$ ; $13$
\end{center}
\begin{questions}[label=\alph*)]
  \item Regrouper ces durées en classes d'amplitude $10$ min, de $\intfo{0}{10}$ à
    $\intfo{40}{50}$, dans un tableau d'effectifs.
  \item Calculer la fréquence de chaque classe, en pourcentage arrondi à $0{,}1\,\%$.
  \item Combien d'élèves mettent moins de $30$ minutes ?
\end{questions}
\begin{correction}
\cas{a)}
\explique{On range chaque durée dans sa classe : $10$ va dans $\intfo{10}{20}$, $20$
dans $\intfo{20}{30}$ et $30$ dans $\intfo{30}{40}$.}{la borne de droite est exclue}
\begin{center}
\renewcommand{\arraystretch}{1.5}
\begin{tabular}{|l|*{5}{c|}c|}
\hline
Durée (min) & $\intfo{0}{10}$ & $\intfo{10}{20}$ & $\intfo{20}{30}$ & $\intfo{30}{40}$
  & $\intfo{40}{50}$ & Total\\\hline
Effectif & $4$ & $10$ & $5$ & $3$ & $2$ & $24$\\\hline
Fréquence & $16{,}7\,\%$ & $41{,}7\,\%$ & $20{,}8\,\%$ & $12{,}5\,\%$ & $8{,}3\,\%$
  & $100\,\%$\\\hline
\end{tabular}
\end{center}
\explique{Vérification : $4+10+5+3+2=\res{24}$, on n'a oublié aucune durée.}{l'effectif
total est celui de la série}
\cas{b)}
\begin{calculs}
f &= \dfrac{10}{24} &\qquad& \rem{classe $\intfo{10}{20}$}\\
f &\approx 0{,}417 &&\\
f &\approx \res{41{,}7\,\%} && \rem{on multiplie par $100$}
\end{calculs}
\explique{Les autres fréquences, calculées de même, sont dans le tableau.}{la somme
vaut $100\,\%$, aux arrondis près}
\cas{c)}
\begin{calculs}
\text{effectif} &= 4+10+5 &\qquad& \rem{les trois premières classes}\\
\text{effectif} &= \res{19} && \rem{$19$ élèves mettent moins de $30$ min}
\end{calculs}
\end{correction}

\etape{À vous de jouer}
Dans une autre classe, $20$ élèves ont indiqué leur temps d'écran de la veille, en
dehors du travail scolaire, en minutes.
\begin{center}
$15$ ; $45$ ; $90$ ; $120$ ; $30$ ; $75$ ; $60$ ; $100$ ; $85$ ; $50$ ; $130$ ; $65$ ;
$40$ ; $95$ ; $110$ ; $70$ ; $20$ ; $60$ ; $140$ ; $80$
\end{center}
\begin{questions}[label=\alph*)]
  \item Regrouper ces durées en classes d'amplitude $30$ min, de $\intfo{0}{30}$ à
    $\intfo{120}{150}$, dans un tableau d'effectifs.
  \item Calculer la fréquence de chaque classe, en pourcentage.
  \item Combien d'élèves ont passé moins de $90$ minutes devant un écran ?
\end{questions}
\reponse{a) effectifs $2$ ; $4$ ; $7$ ; $4$ ; $3$ ; b)~$10\,\%$ ; $20\,\%$ ;
$35\,\%$ ; $20\,\%$ ; $15\,\%$ ; c)~$13$ élèves.}
\end{methodebox}

\afaire{exercice 5}

\partiecours{Moyenne, moyenne pondérée}

\begin{definitionbox}[Définition 5 --- Moyenne]
\begin{itemize}
  \item Moyenne d'une liste de $N$ valeurs :
    $\overline{x}=\dfrac{\text{somme des valeurs}}{N}$.
  \item Valeurs $x_1$, $x_2$, \ldots, $x_p$ d'effectifs $n_1$, $n_2$, \ldots, $n_p$
    (moyenne \textbf{pondérée} par les effectifs) :
    \[
    \overline{x}=\frac{n_1x_1+n_2x_2+\dots+n_px_p}{N}
    \]
  \item Avec les fréquences $f_1$, $f_2$, \ldots, $f_p$ :
    $\overline{x}=f_1x_1+f_2x_2+\dots+f_px_p$.
\end{itemize}
\end{definitionbox}

\begin{methodebox}[Méthode 3 --- Calculer une moyenne]
\etapeexemples
\begin{questions}[label=\alph*)]
  \item Une élève a noté ses heures de sommeil pendant une semaine :
    \begin{center}
    $8$ ; $7{,}5$ ; $9$ ; $6{,}5$ ; $8$ ; $10$ ; $7$
    \end{center}
    Calculer la durée moyenne d'une nuit.
  \item Dans une classe, $28$ élèves ont indiqué leur nombre de séances de sport par
    semaine, hors EPS. Calculer le nombre moyen de séances.
    \begin{center}
    \renewcommand{\arraystretch}{1.5}
    \begin{tabular}{|l|*{5}{c|}}
    \hline
    Nombre de séances & $0$ & $1$ & $2$ & $3$ & $4$\\\hline
    Effectif & $7$ & $6$ & $8$ & $4$ & $3$\\\hline
    \end{tabular}
    \end{center}
  \item Le tableau répartit les ménages français selon leur nombre de personnes en
    $2023$ (chiffres Insee). Calculer le nombre moyen de personnes par ménage, en
    comptant $6$ pour « $6$ ou plus ».
    \begin{center}
    \renewcommand{\arraystretch}{1.5}
    \begin{tabular}{|l|*{6}{c|}}
    \hline
    Nombre de personnes & $1$ & $2$ & $3$ & $4$ & $5$ & $6$ ou plus\\\hline
    Fréquence & $38{,}9\,\%$ & $32{,}3\,\%$ & $12{,}9\,\%$ & $10{,}6\,\%$ & $3{,}8\,\%$
      & $1{,}6\,\%$\\\hline
    \end{tabular}
    \end{center}
\end{questions}
\begin{correction}
\cas{a)}
\begin{calculs}
\overline{x} &= \dfrac{8+7{,}5+9+6{,}5+8+10+7}{7} &\qquad& \rem{$7$ nuits}\\
\overline{x} &= \dfrac{56}{7} &&\\
\overline{x} &= \res{8} && \rem{$8$ heures par nuit en moyenne}
\end{calculs}
\cas{b)}
\begin{calculs}
\overline{x} &= \dfrac{7\times 0+6\times 1+8\times 2+4\times 3+3\times 4}{28}
  &\qquad& \rem{chaque valeur compte autant de fois que son effectif}\\
\overline{x} &= \dfrac{0+6+16+12+12}{28} &&\\
\overline{x} &= \dfrac{46}{28} &&\\
\overline{x} &\approx \res{1{,}64} && \rem{environ $1{,}6$ séance par semaine}
\end{calculs}
\cas{c)}
\begin{calculs}
\overline{x} &= 0{,}389\times 1+0{,}323\times 2+0{,}129\times 3+0{,}106\times 4
  +0{,}038\times 5+0{,}016\times 6 &\qquad& \rem{fréquences en écriture décimale}\\
\overline{x} &= 0{,}389+0{,}646+0{,}387+0{,}424+0{,}19+0{,}096 &&\\
\overline{x} &\approx \res{2{,}13} &&
\end{calculs}
\explique{Un ménage compte en moyenne environ $2{,}13$ personnes ; l'Insee annonce
$2{,}14$, car certains ménages ont plus de $6$ personnes.}{on a compté $6$ pour
« $6$ ou plus »}
\end{correction}

\etape{À vous de jouer}
\begin{questions}[label=\alph*)]
  \item Un élève a chronométré son trajet en bus pendant cinq jours, en minutes :
    \begin{center}
    $22$ ; $25$ ; $19$ ; $24$ ; $30$
    \end{center}
    Calculer la durée moyenne d'un trajet.
  \item Dans une classe, $30$ élèves ont indiqué le nombre de livres lus depuis la
    rentrée, hors cours. Calculer le nombre moyen de livres lus.
    \begin{center}
    \renewcommand{\arraystretch}{1.5}
    \begin{tabular}{|l|*{5}{c|}}
    \hline
    Nombre de livres & $0$ & $1$ & $2$ & $3$ & $4$\\\hline
    Effectif & $6$ & $9$ & $7$ & $5$ & $3$\\\hline
    \end{tabular}
    \end{center}
  \item Le tableau répartit les ménages français selon leur nombre de personnes en
    $1990$ (chiffres Insee). Calculer le nombre moyen de personnes par ménage, en
    comptant $6$ pour « $6$ ou plus », puis comparer avec $2023$.
    \begin{center}
    \renewcommand{\arraystretch}{1.5}
    \begin{tabular}{|l|*{6}{c|}}
    \hline
    Nombre de personnes & $1$ & $2$ & $3$ & $4$ & $5$ & $6$ ou plus\\\hline
    Fréquence & $27{,}0\,\%$ & $29{,}4\,\%$ & $17{,}7\,\%$ & $15{,}7\,\%$ & $6{,}8\,\%$
      & $3{,}4\,\%$\\\hline
    \end{tabular}
    \end{center}
\end{questions}
\reponse{a) $24$ min ; b)~$\frac{50}{30}\approx 1{,}67$ livre ; c)~environ $2{,}56$
personnes (l'Insee annonce $2{,}59$) : les ménages sont plus petits en $2023$.}
\end{methodebox}

\afaire{exercices 6 et 7}

\begin{proprietebox}[Propriété 1 --- Linéarité de la moyenne]
\begin{itemize}
  \item Si on ajoute un même nombre $b$ à toutes les valeurs, la moyenne augmente
    de $b$.
  \item Si on multiplie toutes les valeurs par un même nombre $a$, la moyenne est
    multipliée par $a$.
\end{itemize}
\enonce{Exemple : un trajet dure en moyenne $24$ min. Des travaux ajoutent $5$ min à
chaque trajet : la moyenne passe à $24+5$, soit $29$ min. Si chaque durée augmente
de $10\,\%$, elle est multipliée par $1{,}1$ (coefficient multiplicateur, chapitre 7),
et la moyenne aussi : $24\times 1{,}1$ vaut $26{,}4$ min.}
\end{proprietebox}

\begin{definitionbox}[Définition 6 --- Centre d'une classe]
Le \textbf{centre} de la classe $\intfo{a}{b}$ est le nombre $\dfrac{a+b}{2}$, au
milieu de la classe.
\end{definitionbox}

\begin{proprietebox}[Propriété 2 --- Moyenne d'une série regroupée en classes]
\begin{itemize}
  \item Si l'on connaît l'effectif $n_1$, $n_2$\ldots{} et la moyenne $m_1$,
    $m_2$\ldots{} de chaque classe (ou de chaque groupe), la moyenne de la série est
    $\overline{x}=\dfrac{n_1m_1+n_2m_2+\dots}{N}$.
  \item Si l'on ne connaît que les effectifs, on suppose les valeurs réparties
    uniformément dans chaque classe : la moyenne d'une classe est alors son centre.
    On obtient une \textbf{estimation} de la moyenne.
\end{itemize}
\end{proprietebox}

\begin{methodebox}[Méthode 4 --- Calculer la moyenne d'une série regroupée en classes]
\etapeexemples
\begin{questions}[label=\alph*)]
  \item Un club de judo compte $12$ adultes, d'âge moyen $34$ ans, et $18$ jeunes,
    d'âge moyen $13$ ans. Calculer l'âge moyen des membres du club.
  \item Le tableau répartit $24$ élèves selon leur temps de trajet entre leur domicile
    et le lycée. Estimer le temps de trajet moyen, en arrondissant à $0{,}1$ min.
    \begin{center}
    \renewcommand{\arraystretch}{1.5}
    \begin{tabular}{|l|*{5}{c|}}
    \hline
    Durée (min) & $\intfo{0}{10}$ & $\intfo{10}{20}$ & $\intfo{20}{30}$
      & $\intfo{30}{40}$ & $\intfo{40}{50}$\\\hline
    Effectif & $4$ & $10$ & $5$ & $3$ & $2$\\\hline
    \end{tabular}
    \end{center}
\end{questions}
\begin{correction}
\cas{a)}
\begin{calculs}
\overline{x} &= \dfrac{12\times 34+18\times 13}{12+18} &\qquad& \rem{propriété 2 :
chaque moyenne pèse son effectif}\\
\overline{x} &= \dfrac{408+234}{30} &&\\
\overline{x} &= \dfrac{642}{30} &&\\
\overline{x} &= \res{21{,}4} && \rem{$21{,}4$ ans}
\end{calculs}
\explique{Attention : $\dfrac{34+13}{2}$, soit $23{,}5$, est faux, car il y a plus de
jeunes que d'adultes.}{les deux groupes n'ont pas le même effectif}
\cas{b)}
\explique{Centres des classes : $\dfrac{0+10}{2}=5$, puis $15$, $25$, $35$ et
$45$.}{définition 6}
\begin{calculs}
\overline{x} &\approx \dfrac{4\times 5+10\times 15+5\times 25+3\times 35+2\times 45}{24}
  &\qquad& \rem{chaque classe est remplacée par son centre}\\
\overline{x} &\approx \dfrac{20+150+125+105+90}{24} &&\\
\overline{x} &\approx \dfrac{490}{24} &&\\
\overline{x} &\approx \res{20{,}4} && \rem{environ $20{,}4$ min}
\end{calculs}
\explique{Avec les $24$ durées exactes, on trouverait $19{,}7$ min : l'estimation en
est proche.}{on a supposé les durées réparties uniformément}
\end{correction}

\etape{À vous de jouer}
\begin{questions}[label=\alph*)]
  \item Une chorale compte $15$ adultes, d'âge moyen $42$ ans, et $25$ jeunes, d'âge
    moyen $15$ ans. Calculer l'âge moyen des choristes.
  \item Le tableau répartit $20$ élèves selon leur temps d'écran de la veille. Estimer
    le temps d'écran moyen.
    \begin{center}
    \renewcommand{\arraystretch}{1.5}
    \begin{tabular}{|l|*{5}{c|}}
    \hline
    Durée (min) & $\intfo{0}{30}$ & $\intfo{30}{60}$ & $\intfo{60}{90}$
      & $\intfo{90}{120}$ & $\intfo{120}{150}$\\\hline
    Effectif & $2$ & $4$ & $7$ & $4$ & $3$\\\hline
    \end{tabular}
    \end{center}
\end{questions}
\reponse{a) $\frac{1005}{40}\approx 25{,}1$ ans ; b)~environ $78$ min.}
\end{methodebox}

\afaire{exercices 8 à 11}

\partiecours{Étendue, médiane, quartiles}

\begin{definitionbox}[Définition 7 --- Étendue, médiane, quartiles]
On range les $N$ valeurs de la série dans l'ordre croissant.
\begin{itemize}
  \item \textbf{Étendue} : plus grande valeur $-$ plus petite valeur.
  \item La \textbf{médiane} $Me$ partage la série en deux groupes de même effectif :
    si $N$ est impair, c'est la valeur du milieu, de rang $\dfrac{N+1}{2}$ ; si $N$ est
    pair, c'est la moyenne des valeurs de rangs $\dfrac{N}{2}$ et $\dfrac{N}{2}+1$.
  \item Le \textbf{premier quartile} $Q_1$ est la plus petite valeur de la série telle
    qu'au moins $25\,\%$ des valeurs lui soient inférieures ou égales : son rang est
    $\dfrac{N}{4}$, arrondi à l'entier supérieur.
  \item Le \textbf{troisième quartile} $Q_3$ est la plus petite valeur de la série
    telle qu'au moins $75\,\%$ des valeurs lui soient inférieures ou égales : son rang
    est $\dfrac{3N}{4}$, arrondi à l'entier supérieur.
  \item \textbf{Écart interquartile} : $Q_3-Q_1$. L'intervalle $\intff{Q_1}{Q_3}$
    contient au moins la moitié des valeurs.
\end{itemize}
\end{definitionbox}

\begin{methodebox}[Méthode 5 --- Déterminer l'étendue, la médiane et les quartiles]
\etapeexemples
\begin{questions}[label=\alph*)]
  \item Nombre de naissances en France, en milliers, de $2014$ à $2025$ (chiffres
    Insee, $2023$ à $2025$ provisoires) :
    \begin{center}
    $819$ ; $799$ ; $784$ ; $770$ ; $759$ ; $753$ ; $735$ ; $742$ ; $726$ ; $678$ ;
    $661$ ; $645$
    \end{center}
    Déterminer l'étendue, la médiane, $Q_1$ et $Q_3$, puis interpréter $Q_1$.
  \item Dans une classe, $28$ élèves ont indiqué leur nombre de séances de sport par
    semaine, hors EPS. Déterminer l'étendue, la médiane, $Q_1$ et $Q_3$, puis
    interpréter la médiane.
    \begin{center}
    \renewcommand{\arraystretch}{1.5}
    \begin{tabular}{|l|*{5}{c|}}
    \hline
    Nombre de séances & $0$ & $1$ & $2$ & $3$ & $4$\\\hline
    Effectif & $7$ & $6$ & $8$ & $4$ & $3$\\\hline
    Effectif cumulé croissant & $7$ & $13$ & $21$ & $25$ & $28$\\\hline
    \end{tabular}
    \end{center}
\end{questions}
\begin{correction}
\cas{a)}
\explique{Série rangée : $645$ ; $661$ ; $678$ ; $726$ ; $735$ ; $742$ ; $753$ ; $759$ ;
$770$ ; $784$ ; $799$ ; $819$.}{on range d'abord, $N=12$}
\begin{calculs}
\text{étendue} &= 819-645 &\qquad& \rem{plus grande $-$ plus petite}\\
\text{étendue} &= \res{174} && \rem{$174$ milliers de naissances}
\end{calculs}
\explique{$N=12$ est pair : la médiane est la moyenne des valeurs de rangs $6$ et
$7$, soit $742$ et $753$.}{$\dfrac{12}{2}=6$}
\begin{calculs}
Me &= \dfrac{742+753}{2} &\qquad&\\
Me &= \res{747{,}5} &&
\end{calculs}
\explique{$\dfrac{12}{4}=3$ : $Q_1$ est la $3^{\text{e}}$ valeur, $Q_1=\res{678}$.}{rang
$\dfrac{N}{4}$}
\explique{$\dfrac{3\times 12}{4}=9$ : $Q_3$ est la $9^{\text{e}}$ valeur,
$Q_3=\res{770}$.}{rang $\dfrac{3N}{4}$}
\explique{Au moins un quart des années, de $2014$ à $2025$, ont compté au plus
$678~000$ naissances.}{interprétation de $Q_1$}
\cas{b)}
\begin{calculs}
\text{étendue} &= 4-0 &\qquad&\\
\text{étendue} &= \res{4} &&
\end{calculs}
\explique{$N=28$ : la médiane est la moyenne des valeurs de rangs $14$ et $15$. Les
rangs $14$ à $21$ ont la valeur $2$ (effectifs cumulés $13$ puis $21$), donc
$Me=\res{2}$.}{on lit les rangs dans les effectifs cumulés}
\explique{$\dfrac{28}{4}=7$ : la $7^{\text{e}}$ valeur est $0$ (rangs $1$ à $7$), donc
$Q_1=\res{0}$.}{rang $\dfrac{N}{4}$}
\explique{$\dfrac{3\times 28}{4}=21$ : la $21^{\text{e}}$ valeur est $2$, donc
$Q_3=\res{2}$.}{rang $\dfrac{3N}{4}$}
\explique{Au moins la moitié des élèves font au plus $2$ séances, et au moins la
moitié en font au moins $2$.}{interprétation de la médiane}
\end{correction}

\etape{À vous de jouer}
\begin{questions}[label=\alph*)]
  \item Taux de natalité en France (naissances pour $1~000$ habitants), de $2014$ à
    $2025$ (chiffres Insee) :
    \begin{center}
    $12{,}4$ ; $12{,}0$ ; $11{,}8$ ; $11{,}5$ ; $11{,}3$ ; $11{,}2$ ; $10{,}9$ ;
    $10{,}9$ ; $10{,}6$ ; $9{,}9$ ; $9{,}6$ ; $9{,}4$
    \end{center}
    Déterminer l'étendue, la médiane, $Q_1$ et $Q_3$, puis interpréter $Q_1$.
  \item Dans une classe, $30$ élèves ont indiqué le nombre de livres lus depuis la
    rentrée, hors cours. Déterminer l'étendue, la médiane, $Q_1$ et $Q_3$, puis
    interpréter la médiane.
    \begin{center}
    \renewcommand{\arraystretch}{1.5}
    \begin{tabular}{|l|*{5}{c|}}
    \hline
    Nombre de livres & $0$ & $1$ & $2$ & $3$ & $4$\\\hline
    Effectif & $6$ & $9$ & $7$ & $5$ & $3$\\\hline
    Effectif cumulé croissant & $6$ & $15$ & $22$ & $27$ & $30$\\\hline
    \end{tabular}
    \end{center}
\end{questions}
\reponse{a) étendue $3$ ; $Me=11{,}05$ ; $Q_1=9{,}9$ ; $Q_3=11{,}5$ ; b)~étendue $4$ ;
$Me=1{,}5$ ; $Q_1=1$ ; $Q_3=3$.}
\end{methodebox}

\afaire{exercices 12 à 14}

\begin{remarquebox}[Effet d'une valeur extrême]
La série $8$ ; $9$ ; $10$ ; $11$ ; $12$ a pour moyenne $10$ et pour médiane $10$. Si
l'on ajoute la valeur $40$, la moyenne devient $\dfrac{90}{6}$, soit $15$, mais la
médiane devient $\dfrac{10+11}{2}$, soit $10{,}5$ seulement. La moyenne est très
sensible à l'ajout ou à la suppression d'une valeur extrême ; la médiane, peu.
\end{remarquebox}

\begin{definitionbox}[Définition 8 --- Classe médiane]
Pour une série regroupée en classes, la \textbf{classe médiane} est la première
classe dont l'effectif cumulé croissant atteint ou dépasse $\dfrac{N}{2}$. Si l'on
suppose les valeurs réparties uniformément dans cette classe, on estime la médiane
par proportionnalité.
\end{definitionbox}

\begin{methodebox}[Méthode 6 --- Estimer la médiane d'une série regroupée en classes]
\etapeexemples
Le tableau répartit $24$ élèves selon leur temps de trajet entre leur domicile et le
lycée.
\begin{center}
\renewcommand{\arraystretch}{1.5}
\begin{tabular}{|l|*{5}{c|}}
\hline
Durée (min) & $\intfo{0}{10}$ & $\intfo{10}{20}$ & $\intfo{20}{30}$ & $\intfo{30}{40}$
  & $\intfo{40}{50}$\\\hline
Effectif & $4$ & $10$ & $5$ & $3$ & $2$\\\hline
\end{tabular}
\end{center}
\begin{questions}[label=\alph*)]
  \item Calculer les effectifs cumulés croissants.
  \item Déterminer la classe médiane.
  \item Estimer la médiane, en supposant les durées réparties uniformément dans
    cette classe.
\end{questions}
\begin{correction}
\cas{a)}
\explique{Effectifs cumulés : $4$ ; $4+10=14$ ; $14+5=19$ ; $19+3=22$ ;
\res{$22+2=24$}.}{chaque effectif s'ajoute au précédent}
\cas{b)}
\explique{$\dfrac{24}{2}=12$ : l'effectif cumulé atteint $12$ dans la classe
$\intfo{10}{20}$ ($4<12\le 14$), qui est la \res{classe médiane}.}{définition 8}
\cas{c)}
\explique{Il faut $12-4=8$ élèves de cette classe, sur ses $10$ élèves : on avance de
$\dfrac{8}{10}$ de l'amplitude de la classe.}{les $4$ premiers sont dans
$\intfo{0}{10}$}
\begin{calculs}
Me &\approx 10+\dfrac{8}{10}\times 10 &\qquad& \rem{début de la classe $+$ part de son
amplitude}\\
Me &\approx \res{18} && \rem{environ $18$ min}
\end{calculs}
\explique{Avec les $24$ durées exactes, on trouverait $17{,}5$ min.}{l'estimation est
proche}
\end{correction}

\etape{À vous de jouer}
Le tableau répartit $20$ élèves selon leur temps d'écran de la veille.
\begin{center}
\renewcommand{\arraystretch}{1.5}
\begin{tabular}{|l|*{5}{c|}}
\hline
Durée (min) & $\intfo{0}{30}$ & $\intfo{30}{60}$ & $\intfo{60}{90}$ & $\intfo{90}{120}$
  & $\intfo{120}{150}$\\\hline
Effectif & $2$ & $4$ & $7$ & $4$ & $3$\\\hline
\end{tabular}
\end{center}
\begin{questions}[label=\alph*)]
  \item Calculer les effectifs cumulés croissants.
  \item Déterminer la classe médiane.
  \item Estimer la médiane, en supposant les durées réparties uniformément dans
    cette classe, au dixième de minute.
\end{questions}
\reponse{a) $2$ ; $6$ ; $13$ ; $17$ ; $20$ ; b)~$\intfo{60}{90}$ ;
c)~$Me\approx 60+\frac{4}{7}\times 30\approx 77{,}1$ min.}
\end{methodebox}

\afaire{exercices 15 et 16}

\partiecours{Écart type}

Deux séries peuvent avoir la même moyenne, mais des valeurs plus ou moins
dispersées autour d'elle. L'écart type mesure cette dispersion.

\begin{definitionbox}[Définition 9 --- Écart type]
L'\textbf{écart type} d'une série de moyenne $\overline{x}$, de valeurs $x_1$, \ldots,
$x_p$ et d'effectifs $n_1$, \ldots, $n_p$, est le nombre
\[
\sigma=\sqrt{\frac{n_1(x_1-\overline{x})^{2}+n_2(x_2-\overline{x})^{2}+\dots
+n_p(x_p-\overline{x})^{2}}{N}}
\]
C'est la racine carrée de la moyenne des carrés des écarts à la moyenne. Il a la même
unité que les valeurs ; plus il est grand, plus les valeurs sont dispersées autour de
la moyenne.
\end{definitionbox}

\begin{remarquebox}[Résumer une série]
On résume une série par un couple d'indicateurs : moyenne et écart type, ou médiane
et écart interquartile. Le second couple est peu sensible aux valeurs extrêmes.
\end{remarquebox}

\begin{proprietebox}[Calculatrice NumWorks --- application \texttt{Statistiques}]
\begin{center}
\renewcommand{\arraystretch}{1.3}
\begin{tabular}{|>{\raggedright\arraybackslash}p{4.3cm}|>{\raggedright\arraybackslash}p{12.3cm}|}
\hline
\textbf{Pour} & \textbf{On fait}\\\hline
entrer la série & Onglet \texttt{Données} : les valeurs dans la colonne \texttt{V1},
  leurs effectifs dans \texttt{N1} (effectif $1$ pour chaque valeur d'une
  liste).\\\hline
lire les indicateurs & Onglet \texttt{Stats} : effectif total, minimum, maximum,
  étendue, moyenne, \texttt{Écart type} $\sigma$, $Q_1$, médiane, $Q_3$, écart
  interquartile.\\\hline
\end{tabular}
\end{center}
On lit l'écart type $\sigma$, pas l'écart type « échantillon » $s$.
{\footnotesize\itshape Les intitulés peuvent varier selon la version du logiciel.}
\end{proprietebox}

\begin{methodebox}[Méthode 7 --- Calculer un écart type et comparer deux séries]
\etapeexemples
Hugo vient au lycée à vélo, Inès en bus. Voici la durée de leurs trajets, en minutes,
pendant une semaine.
\begin{center}
Hugo : $14$ ; $15$ ; $16$ ; $15$ ; $15$ \qquad Inès : $10$ ; $25$ ; $12$ ; $18$ ; $10$
\end{center}
\begin{questions}[label=\alph*)]
  \item Calculer à la main la moyenne et l'écart type des durées de Hugo.
  \item Déterminer à la calculatrice la moyenne et l'écart type des durées d'Inès,
    arrondi au centième.
  \item Comparer les deux séries.
\end{questions}
\begin{correction}
\cas{a)}
\begin{calculs}
\overline{x} &= \dfrac{14+15+16+15+15}{5} &\qquad&\\
\overline{x} &= \dfrac{75}{5} &&\\
\overline{x} &= \res{15} &&
\end{calculs}
\explique{On calcule l'écart de chaque valeur à la moyenne, puis son carré.}{un carré
n'est jamais négatif}
\begin{center}
\renewcommand{\arraystretch}{1.5}
\begin{tabular}{|l|*{5}{c|}}
\hline
Durée $x$ & $14$ & $15$ & $16$ & $15$ & $15$\\\hline
Écart $x-15$ & $-1$ & $0$ & $1$ & $0$ & $0$\\\hline
Carré de l'écart & $1$ & $0$ & $1$ & $0$ & $0$\\\hline
\end{tabular}
\end{center}
\begin{calculs}
\sigma &= \sqrt{\dfrac{1+0+1+0+0}{5}} &\qquad& \rem{moyenne des carrés, puis racine}\\
\sigma &= \sqrt{0{,}4} &&\\
\sigma &\approx \res{0{,}63} && \rem{environ $0{,}63$ min}
\end{calculs}
\cas{b)}
\explique{Les cinq durées dans \texttt{V1}, $1$ dans chaque case de \texttt{N1} ;
l'onglet \texttt{Stats} donne $\overline{x}=\res{15}$ et
$\sigma\approx\res{5{,}80}$.}{application \texttt{Statistiques}}
\cas{c)}
\explique{Les deux élèves mettent en moyenne $15$ min, mais l'écart type d'Inès est
environ $9$ fois plus grand : ses trajets sont \res{beaucoup plus irréguliers} que
ceux de Hugo.}{même moyenne, dispersions très différentes}
\end{correction}

\etape{À vous de jouer}
Léa et Tom ont fait cinq quiz de révision notés sur $20$.
\begin{center}
Léa : $8$ ; $10$ ; $12$ ; $10$ ; $10$ \qquad Tom : $4$ ; $16$ ; $6$ ; $14$ ; $10$
\end{center}
\begin{questions}[label=\alph*)]
  \item Calculer à la main la moyenne et l'écart type des notes de Léa.
  \item Déterminer à la calculatrice la moyenne et l'écart type des notes de Tom,
    arrondi au centième.
  \item Comparer les deux séries.
\end{questions}
\reponse{a) $\overline{x}=10$ ; $\sigma=\sqrt{1{,}6}\approx 1{,}26$ ; b)~$\overline{x}=10$ ;
$\sigma\approx 4{,}56$ ; c)~même moyenne, mais les notes de Tom sont bien plus
dispersées.}
\end{methodebox}

\afaire{exercices 17 à 20}

\partiecours{Diagrammes}

\begin{definitionbox}[Définition 10 --- Choisir un diagramme]
\begin{itemize}
  \item Caractère qualitatif : \textbf{diagramme circulaire} (l'angle de chaque
    secteur est proportionnel à l'effectif : angle $=$ fréquence $\times 360^{\circ}$)
    ou \textbf{diagramme en barres}.
  \item Caractère quantitatif discret : diagramme en barres, une barre par valeur.
  \item Série regroupée en classes de même amplitude : \textbf{histogramme}, des
    rectangles accolés, de hauteur proportionnelle à l'effectif de la classe.
  \item \textbf{Polygone des fréquences cumulées croissantes} : on place le point
    (début de la première classe ; $0$), puis pour chaque classe le point (fin de la
    classe ; fréquence cumulée), et on les relie par des segments. Le point
    d'ordonnée $50\,\%$ donne une estimation de la médiane.
\end{itemize}
\end{definitionbox}

\begin{methodebox}[Méthode 8 --- Construire un diagramme circulaire]
\etapeexemples
Répartition de la population française au $1^{\text{er}}$ janvier $2026$ selon l'âge
(chiffres Insee, provisoires).
\begin{center}
\renewcommand{\arraystretch}{1.5}
\begin{tabular}{|l|c|c|c|}
\hline
Âge & moins de $20$ ans & de $20$ à $64$ ans & $65$ ans ou plus\\\hline
Fréquence & $22{,}5\,\%$ & $55{,}3\,\%$ & $22{,}2\,\%$\\\hline
\end{tabular}
\end{center}
\begin{questions}[label=\alph*)]
  \item Calculer l'angle de chaque secteur, arrondi au degré.
  \item Construire le diagramme circulaire.
\end{questions}
\begin{correction}
\cas{a)}
\begin{calculs}
\text{angle} &= 0{,}225\times 360 &\qquad& \rem{moins de $20$ ans}\\
\text{angle} &= \res{81^{\circ}} &&
\end{calculs}
\begin{calculs}
\text{angle} &= 0{,}553\times 360 &\qquad& \rem{de $20$ à $64$ ans}\\
\text{angle} &\approx \res{199^{\circ}} && \rem{$199{,}08$ arrondi}
\end{calculs}
\begin{calculs}
\text{angle} &= 0{,}222\times 360 &\qquad& \rem{$65$ ans ou plus}\\
\text{angle} &\approx \res{80^{\circ}} && \rem{$79{,}92$ arrondi}
\end{calculs}
\explique{Vérification : $81+199+80=360$.}{les secteurs font le tour complet}
\cas{b)}
\par\nopagebreak
\begin{minipage}[c]{0.58\linewidth}\raggedright
On trace un cercle, un premier rayon, puis chaque angle au rapporteur, à la suite du
précédent ; on écrit dans chaque secteur ce qu'il représente.
\end{minipage}\hfill
\begin{minipage}[c]{0.38\linewidth}\centering
\begin{tikzpicture}
  \fill[coulCourbe!20] (0,0) -- (90:2) arc (90:171:2) -- cycle;
  \fill[coulCourbeC!20] (0,0) -- (171:2) arc (171:370:2) -- cycle;
  \fill[coulCourbeB!20] (0,0) -- (10:2) arc (10:90:2) -- cycle;
  \draw[secteur] (0,0) circle (2);
  \foreach \a in {90,171,10} \draw[secteur] (0,0) -- (\a:2);
  \node[font=\footnotesize,align=center] at (131:1.25) {moins\\de 20 ans};
  \node[font=\footnotesize,align=center] at (270:1) {20 à 64 ans};
  \node[font=\footnotesize,align=center] at (50:1.25) {65 ans\\ou plus};
\end{tikzpicture}
\end{minipage}
\end{correction}

\etape{À vous de jouer}
Mode de transport principal des actifs pour aller travailler en France (chiffres
Insee, recensement $2017$).
\begin{center}
\renewcommand{\arraystretch}{1.5}
\begin{tabular}{|l|*{5}{c|}}
\hline
Mode & voiture & transports en commun & marche & vélo & deux-roues motorisé\\\hline
Fréquence & $73{,}7\,\%$ & $15{,}8\,\%$ & $6{,}3\,\%$ & $2{,}2\,\%$ & $2{,}0\,\%$\\\hline
\end{tabular}
\end{center}
\begin{questions}[label=\alph*)]
  \item Calculer l'angle de chaque secteur, arrondi au degré.
  \item Construire le diagramme circulaire.
\end{questions}
\reponse{a) $265^{\circ}$ ; $57^{\circ}$ ; $23^{\circ}$ ; $8^{\circ}$ ; $7^{\circ}$ (total
$360^{\circ}$).}
\end{methodebox}

\afaire{exercices 21 et 22}

\begin{methodebox}[Méthode 9 --- Construire un histogramme et le polygone des fréquences cumulées]
\etapeexemples
Le tableau répartit $24$ élèves selon leur temps de trajet entre leur domicile et le
lycée.
\begin{center}
\renewcommand{\arraystretch}{1.5}
\begin{tabular}{|l|*{5}{c|}}
\hline
Durée (min) & $\intfo{0}{10}$ & $\intfo{10}{20}$ & $\intfo{20}{30}$ & $\intfo{30}{40}$
  & $\intfo{40}{50}$\\\hline
Effectif & $4$ & $10$ & $5$ & $3$ & $2$\\\hline
\end{tabular}
\end{center}
\begin{questions}[label=\alph*)]
  \item Construire l'histogramme de cette série.
  \item Calculer les fréquences cumulées croissantes, en pourcentage arrondi à
    $0{,}1\,\%$, puis construire le polygone des fréquences cumulées.
  \item Lire sur le polygone une estimation de la médiane.
\end{questions}
\begin{correction}
\cas{a)}
\par\nopagebreak
\begin{minipage}[c]{0.5\linewidth}\raggedright
Axe horizontal : la durée, $1$ carreau pour $5$ min ; axe vertical : l'effectif,
$1$ carreau pour $1$ élève. Chaque classe a pour rectangle sa largeur ($10$ min)
et une hauteur égale à son effectif ; les rectangles sont accolés.
\end{minipage}\hfill
\begin{minipage}[c]{0.46\linewidth}\centering
\begin{tikzpicture}[x=0.5cm,y=0.4cm]
  \quadrillage{0}{0}{11}{11}
  \foreach \a/\h in {0/4,2/10,4/5,6/3,8/2} \draw[barre] (\a,0) rectangle ({\a+2},\h);
  \draw[-{Stealth[length=2mm]},thick] (0,0)--(11.5,0) node[above left,font=\footnotesize]{min};
  \draw[-{Stealth[length=2mm]},thick] (0,0)--(0,11.5) node[left,font=\footnotesize]{effectif};
  \foreach \k/\l in {0/0,2/10,4/20,6/30,8/40,10/50}
    \draw (\k,0.15)--(\k,-0.15) node[below,font=\scriptsize]{$\l$};
  \foreach \k in {2,4,6,8,10}
    \draw (0.12,\k)--(-0.12,\k) node[left,font=\scriptsize]{$\k$};
\end{tikzpicture}
\end{minipage}
\cas{b)}
\explique{Fréquences cumulées : $\dfrac{4}{24}\approx 16{,}7\,\%$, puis
$\dfrac{14}{24}\approx 58{,}3\,\%$, $\dfrac{19}{24}\approx 79{,}2\,\%$,
$\dfrac{22}{24}\approx 91{,}7\,\%$ et \res{$100\,\%$}.}{effectifs cumulés $4$ ; $14$ ;
$19$ ; $22$ ; $24$}
\par\nopagebreak
\begin{minipage}[c]{0.5\linewidth}\raggedright
On place les points $(0\,;0)$, $(10\,;16{,}7)$, $(20\,;58{,}3)$, $(30\,;79{,}2)$,
$(40\,;91{,}7)$ et $(50\,;100)$ : chaque fréquence cumulée est placée à la fin de sa
classe. On les relie par des segments.

\smallskip
\cas{c)} On trace la droite horizontale d'ordonnée $50\,\%$ jusqu'au polygone, puis on
descend sur l'axe : $Me\approx\res{18}$ min, comme à la méthode 6.
\end{minipage}\hfill
\begin{minipage}[c]{0.46\linewidth}\centering
\begin{tikzpicture}[x=0.5cm,y=0.4cm]
  \quadrillage{0}{0}{11}{10}
  \draw[-{Stealth[length=2mm]},thick] (0,0)--(11.5,0) node[above left,font=\footnotesize]{min};
  \draw[-{Stealth[length=2mm]},thick] (0,0)--(0,10.6) node[right,font=\footnotesize]{\%};
  \foreach \k/\l in {0/0,2/10,4/20,6/30,8/40,10/50}
    \draw (\k,0.15)--(\k,-0.15) node[below,font=\scriptsize]{$\l$};
  \foreach \k/\l in {2/20,5/50,8/80,10/100}
    \draw (0.12,\k)--(-0.12,\k) node[left,font=\scriptsize]{$\l$};
  \draw[red,dashed,thick] (0,5)--(3.6,5)--(3.6,0);
  \draw[coulCourbe,line width=1.1pt] (0,0)--(2,1.67)--(4,5.83)--(6,7.92)--(8,9.17)--(10,10);
  \foreach \p in {(0,0),(2,1.67),(4,5.83),(6,7.92),(8,9.17),(10,10)}
    \path[coulCourbe] \p node[croix]{};
\end{tikzpicture}
\end{minipage}
\end{correction}

\etape{À vous de jouer}
Le tableau répartit $20$ élèves selon leur temps d'écran de la veille.
\begin{center}
\renewcommand{\arraystretch}{1.5}
\begin{tabular}{|l|*{5}{c|}}
\hline
Durée (min) & $\intfo{0}{30}$ & $\intfo{30}{60}$ & $\intfo{60}{90}$ & $\intfo{90}{120}$
  & $\intfo{120}{150}$\\\hline
Effectif & $2$ & $4$ & $7$ & $4$ & $3$\\\hline
\end{tabular}
\end{center}
\begin{questions}[label=\alph*)]
  \item Construire l'histogramme de cette série.
  \item Calculer les fréquences cumulées croissantes, en pourcentage, puis construire
    le polygone des fréquences cumulées.
  \item Lire sur le polygone une estimation de la médiane.
\end{questions}
\reponse{b) $10\,\%$ ; $30\,\%$ ; $65\,\%$ ; $85\,\%$ ; $100\,\%$ ; c)~environ
$77$ min.}
\end{methodebox}

\afaire{exercices 23 et 24}

\partiecours{Croisement de deux variables qualitatives}

Un tableau croisé (chapitre 11) répartit une population selon deux caractères. On y
lit deux sortes de fréquences.

\begin{definitionbox}[Définition 11 --- Fréquences marginales et conditionnelles]
\begin{itemize}
  \item La \textbf{fréquence marginale} d'une valeur se calcule avec un total, dans
    la marge du tableau :
    \[
    f=\frac{\text{total de sa ligne (ou de sa colonne)}}{\text{effectif total}}
    \]
  \item La \textbf{fréquence conditionnelle} de $B$ \textbf{parmi} $A$ est la
    proportion des individus qui ont $B$ parmi ceux qui ont $A$ :
    \[
    f=\frac{\text{effectif de } A \text{ et } B}{\text{effectif de } A}
    \]
    On ne regarde que la ligne (ou la colonne) de $A$.
\end{itemize}
Pour un individu choisi au hasard, c'est la probabilité $P_A(B)$ du chapitre 11. Les
médias l'écrivent « parmi les… », « chez les… », « des … , $x\,\%$ … ».
\end{definitionbox}

\begin{methodebox}[Méthode 10 --- Calculer des fréquences marginales et conditionnelles]
\etapeexemples
Le tableau répartit les actifs selon leur mode de transport principal pour aller
travailler et la distance entre leur domicile et leur travail, en milliers d'actifs
(chiffres Insee, recensement $2017$, arrondis à la dizaine de milliers).
\begin{center}
\renewcommand{\arraystretch}{1.5}
\begin{tabular}{|l|c|c|c|c|}
\hline
 & $2$ km ou moins & de $2$ à $5$ km & plus de $5$ km & Total\\\hline
Marche ou vélo & $1~320$ & $540$ & $250$ & $2~110$\\\hline
Voiture & $1~960$ & $2~990$ & & \\\hline
Autre mode & $420$ & $970$ & $2~960$ & $4~350$\\\hline
Total & $3~700$ & $4~500$ & $16~400$ & $24~600$\\\hline
\end{tabular}
\end{center}
\begin{questions}[label=\alph*)]
  \item Parmi les actifs qui habitent à plus de $5$ km de leur travail, $80{,}4\,\%$
    prennent la voiture. Compléter le tableau.
  \item Calculer la fréquence marginale des actifs qui prennent la voiture.
  \item Calculer la fréquence des actifs qui vont travailler à pied ou à vélo parmi
    ceux qui habitent à $2$ km ou moins.
  \item Calculer la fréquence des actifs qui habitent à $2$ km ou moins parmi ceux qui
    vont travailler à pied ou à vélo. Traduire le résultat par une phrase.
\end{questions}
\begin{correction}
\cas{a)}
\begin{calculs}
\text{effectif} &= 0{,}804\times 16~400 &\qquad& \rem{$80{,}4\,\%$ de la colonne
« plus de $5$ km »}\\
\text{effectif} &= 13~185{,}6 &&\\
\text{effectif} &\approx \res{13~190} && \rem{arrondi à la dizaine}
\end{calculs}
\begin{calculs}
\text{total voiture} &= 1~960+2~990+13~190 &\qquad& \rem{somme de la ligne}\\
\text{total voiture} &= \res{18~140} &&
\end{calculs}
\explique{Vérification : $250+13~190+2~960=16~400$ et
$2~110+18~140+4~350=24~600$.}{chaque colonne a pour somme son total}
\cas{b)}
\begin{calculs}
f &= \dfrac{18~140}{24~600} &\qquad& \rem{total de la ligne sur l'effectif total}\\
f &\approx \res{0{,}737} && \rem{soit $73{,}7\,\%$ des actifs}
\end{calculs}
\cas{c)}
\begin{calculs}
f &= \dfrac{1~320}{3~700} &\qquad& \rem{on ne regarde que la colonne « $2$ km ou
moins »}\\
f &\approx \res{0{,}357} && \rem{soit environ $35{,}7\,\%$}
\end{calculs}
\cas{d)}
\begin{calculs}
f &= \dfrac{1~320}{2~110} &\qquad& \rem{on ne regarde que la ligne « marche ou vélo »}\\
f &\approx \res{0{,}626} && \rem{soit environ $62{,}6\,\%$}
\end{calculs}
\explique{Environ $63\,\%$ des actifs qui vont travailler à pied ou à vélo habitent à
$2$ km ou moins de leur travail.}{même case qu'en c), autre population de
référence}
\end{correction}

\etape{À vous de jouer}
Le tableau répartit les $600$ élèves d'un lycée selon leur moyen de transport et la
durée de leur trajet.
\begin{center}
\renewcommand{\arraystretch}{1.5}
\begin{tabular}{|l|c|c|c|c|}
\hline
 & moins de $15$ min & de $15$ à $30$ min & plus de $30$ min & Total\\\hline
Marche ou vélo & $120$ & $45$ & $5$ & $170$\\\hline
Bus ou car & $30$ & & $110$ & \\\hline
Voiture & $50$ & $80$ & $35$ & $165$\\\hline
Total & $200$ & $250$ & $150$ & $600$\\\hline
\end{tabular}
\end{center}
\begin{questions}[label=\alph*)]
  \item Parmi les élèves dont le trajet dure de $15$ à $30$ min, $50\,\%$ prennent le
    bus ou le car. Compléter le tableau.
  \item Calculer la fréquence marginale des élèves qui prennent le bus ou le car.
  \item Calculer la fréquence des élèves qui viennent à pied ou à vélo parmi ceux dont
    le trajet dure moins de $15$ min.
  \item Calculer la fréquence des élèves dont le trajet dure moins de $15$ min parmi
    ceux qui viennent à pied ou à vélo. Traduire le résultat par une phrase.
\end{questions}
\reponse{a) $125$ ; total $265$ ; b)~$\frac{265}{600}\approx 44{,}2\,\%$ ;
c)~$\frac{120}{200}=60\,\%$ ; d)~$\frac{120}{170}\approx 70{,}6\,\%$ des élèves qui
viennent à pied ou à vélo ont un trajet de moins de $15$ min.}
\end{methodebox}

\afaire{exercices 25 à 28}

\end{document}
